% Part: first-order-logic
% Chapter: proof-systems
% Section: seqeunt-calculus

\documentclass[../../../include/open-logic-section]{subfiles}

\begin{document}

\iftag{FOL}
      {\olfileid[hi]{fol}{prf}{seq}}
      {\olfileid[hi]{pl}{prf}{seq}}

\olsection{अनुक्रम-कलन}

जहाँ अनेक !!{derivation}-तंत्र !!{sentence}s की विभिन्न व्यवस्थाओं पर कार्य
करते हैं, वहीं अनुक्रम-कलन \emph{अनुक्रमों} पर कार्य करता है। अनुक्रम इस रूप
का व्यंजक होता है:
\[
!A_1, \dots, !A_m \Sequent !B_1, \dots, !B_m,
\]
अर्थात् यह !!{sentence}s के दो क्रमों का युग्म है, जिन्हें अनुक्रम-चिह्न
$\Sequent$ अलग करता है। दोनों में से कोई भी क्रम रिक्त हो सकता है।
अनुक्रम-कलन में !!^{derivation}, अनुक्रमों का एक वृक्ष होती है। उसके सबसे ऊपर
के अनुक्रम विशेष रूप के होते हैं (उन्हें ``आरंभिक अनुक्रम'' अथवा ``अभिगृहीत''
कहते हैं), और हर अन्य अनुक्रम अपने ठीक ऊपर के अनुक्रमों से किसी अनुमान-नियम
द्वारा निकलता है। अनुमान-नियम या तो अनुक्रमों के !!{sentence}s को बदलते हैं
(बाईं या दाईं ओर उन्हें जोड़कर, हटाकर अथवा पुनर्व्यवस्थित करके), या फिर नियम
के निष्कर्ष में कोई जटिल !!{formula} प्रस्तुत करते हैं। उदाहरण के लिए,
$\LeftR{\land}$ नियम $!A,
\Gamma \Sequent \Delta$ से $!A \land !B, \Gamma \Sequent \Delta$ का अनुमान
लगाने देता है; और $\RightR{\lif}$ नियम $!A, \Gamma \Sequent
\Delta, !B$ से $\Gamma \Sequent \Delta, !A \lif !B$ का अनुमान लगाने देता
है, जहाँ $\Gamma$, $\Delta$, $!A$ और~$!B$ कोई भी हों। (विशेषतः, $\Gamma$
और~$\Delta$ रिक्त हो सकते हैं।)

अनुक्रम-कलन पर आधारित $\Proves$ संबंध की परिभाषा इस प्रकार है:
$\Gamma \Proves !A$ तभी और केवल तभी जब कोई क्रम $\Gamma_0$ ऐसा हो कि हर
$!A$, जो $\Gamma_0$ में है,~$\Gamma$ में हो और कोई ऐसी !!{derivation} हो जिसके
मूल पर अनुक्रम~$\Gamma_0 \Sequent !A$ हो। $!A$ अनुक्रम-कलन में प्रमेय है यदि
अनुक्रम~$\Sequent !A$ की कोई !!{derivation} हो। उदाहरण के लिए, यह
!!{derivation} दिखाती है कि $\Proves (!A \land !B) \lif !A$:
\begin{prooftree}
  \Axiom$!A \fCenter !A$
  \RightLabel{\LeftR{\land}}
  \UnaryInf$!A \land !B \fCenter !A$
  \RightLabel{\RightR{\lif}}
  \UnaryInf$\fCenter (!A \land !B) \lif !A$
\end{prooftree}

अनुक्रम-कलन में समुच्चय $\Gamma$ असंगत है यदि $\Gamma_0 \Sequent$ की कोई
!!{derivation} हो (जहाँ हर $!A \in \Gamma_0$,~$\Gamma$ में हो और अनुक्रम का
दायाँ पक्ष रिक्त हो)। नियम \RightR{\Weakening} का उपयोग करके असंगत समुच्चय से
कोई भी !!{sentence} !!{derive} किया जा सकता है।

अनुक्रम-कलन का आविष्कार गेरहार्ड गेन्ट्सेन ने 1930 के दशक में किया। अपने
व्यवस्थित और सममित विन्यास के कारण यह बहुत उपयोगी औपचारिक ढाँचा है; इसकी
सहायता से यह सिद्धान्त विकसित किया जा सकता है कि !!{derivation}s कैसे काम
करती हैं। अनुक्रम-कलन में !!{derivation}s
खोजना अपेक्षाकृत आसान है, पर ये !!{derivation}s प्रायः पढ़ने में कठिन होती हैं
और प्रमाणों से इनका संबंध भी कभी-कभी आसानी से दिखाई नहीं देता। फिर भी
!!{derivation}-तंत्रों
के प्रति यह अत्यंत सुरुचिपूर्ण दृष्टिकोण सिद्ध हुआ है, और अनेक तर्कशास्त्रों
के अपने अनुक्रम-कलन तंत्र हैं।

\end{document}
